%%%%%%%%%%%%%%%%%%%%%%%%%%%%%%%%%%%%%%%%%%%%%%%%%%%%%
%			HLAVIČKA								%
%%%%%%%%%%%%%%%%%%%%%%%%%%%%%%%%%%%%%%%%%%%%%%%%%%%%%
\documentclass[openany,12pt]{memoir}

\usepackage[utf8]{inputenc} 
\usepackage[czech]{babel}
\usepackage[T1]{fontenc}
\usepackage[top=1.5cm, bottom=2cm, left=2cm, right=2cm]{geometry}  % --> NASTAVENÍ OKRAJŮ
\usepackage{fancyhdr}
\usepackage{graphicx}
\usepackage{lmodern}
\usepackage{xwatermark}
\usepackage{xcolor}
\usepackage{changepage}
\usepackage{newunicodechar}



%%%%%% Package na zpěvník
\usepackage[full]{leadsheets}%http://mirrors.nic.cz/tex-archive/macros/latex/contrib/leadsheets/leadsheets_en.pdf   --> dokumentace	
\definesongtitletemplate{empty}{} 
\setchords{
format = \bfseries,   %tučné akordy
minor = {mi},% 
input-notation = {german},%
output-notation = {german}%
}
\definesongtitletemplate{empty}{} 

\newlength{\drop}
\newwatermark[allpages,color=red!50,angle=0,scale=2, xpos=0,ypos=0]{\includegraphics[width=5cm]{obr/pozadi.jpg}} %--> dvojka na pozadí


%%%% Vlastní příkazy
\newcounter{Slokočet}   %Automatické číslování slok
\newcommand{\mezera}{\vspace*{0.5cm}}   %Horizontální odsazení slok
\newcommand{\stred}{5.2cm}   %%% Na zarovnání slok doprostřed, pozn. automatičtější zarovnávání na střed nejde
\newcommand{\refren}{\mezera \noindent \textbf{R:} } %refrén
\newcommand{\sloka}{\mezera \noindent \addtocounter{Slokočet}{1} \arabic{Slokočet}. } 	%sloka, která se automaticky čísluje
\newcommand{\ssloka}{\mezera \noindent}  % vlastní číslo sloky
\newcommand{\carka}{,\:}
\newcommand{\m}[1]{\color{white}{#1}}  %Pro akordy

%%%%%%%%%%%%%%%%%%%%%%%%%%%%%%%%%%%%%%%%%%%%%%%%%%%%%%
%			NÁVOD									 %
%%%%%%%%%%%%%%%%%%%%%%%%%%%%%%%%%%%%%%%%%%%%%%%%%%%%%%
%1. Věci v hlavičce IGNOROVAT
%2. Píseň psát do prostoru mezi \begin{song} a \end{song}
%3. další řádek se značí dvěma odsazeníma (= dvakrát stisknout enter)
%4. \refren vždy na začátku refrenu a \sloka na začátek sloky (automaticky se čísluje)
%  \ = alt gr + q ; [/] = alt gr f/g ; {/} = alt gr + b/n; ^ = alt gr + 3 + mezera
%Cokoliv napíšete do ^{  } se bude brát jako akord
%když se toto bude dotýkat nějakého slova (nebude mezi tím a slovem mezera)
%tak se akordy zjeví nad slovem, ale pište to před slovo
%Když se to nedotýká slova, tak akord lítá ve vzduchu a vytiskne se větší mezera
%První možnost je asi preferovanější
%5. Akordy stačí psát jen do první sloky, když se nezmění -- kytaristi to zvládnou



%<++>
%\usepackage{subfiles}
%</++>

\begin{document}

\pagestyle{empty}
%\addtocontents{toc}{\protect\thispagestyle{empty}}
%\tableofcontents \thispagestyle{empty}\newpage
\newgeometry{top=1.5cm, bottom = 0cm, left = 2cm, right = 2cm}




%Pro kompilaci je důležitý, aby ten input a \end{document} zůstali na posledních
%dvou řádcích!!!!
\documentclass[../main.tex]{subfiles}

\chapter*{Do Ameriky}

\noindent\hspace{0.15\linewidth}
\begin{minipage}{0.7\linewidth}
\begin{song}{title=Do Ameriky}

\begin{verse}[numbered]
	^{D}Do Ameriky jezděj Parníky,\\
	když tam přijdeš, zdá se ti to ^{A7}všecko veliký,\\
	^{\phantom{S}}je to fakticky tuze praktický  \\
	když tam přijdeš, umět anglic^{D}ky. \\
\end{verse}

\begin{verse*}
	Refrain:
	^{D}Dobrou noc - good night, ^{A7}výborně - all right,\\
	Konrád White, už je off^{D}side (housemister - Dally sister!) \\
	^{D}Dneska je today, ^{A7}dobrý den - good day,\\
	já jsem byl taky chu^{D}dej. \\ 
	^{G}His Master's Voice, Yankee Doo^{G}dle,\\ \\
	^{En7}máš apetit? (mám) ^{A7}vem si štrůdl! (do pusy)\\
	^{D}Sto kilo je cent, ^{A7}patent je patent,\\
	husa v troubě happy^{D}end.
\end{verse*}

\begin{verse}[numbered]
	Co Buffalo Bill vrazil do kobyl,\\
	za to by si bejval byl koupil automobil,\\
	cowboy na koni už se nehoní,\\
	přestože mu Fordka nevoní.\\
\end{verse}

\begin{verse*}
	Refrain
\end{verse*}

\begin{verse}[numbered]
	Mám na západě chajdu v Nevadě,\\
	zlatou žílu v Arizoně, doly v Kanadě,\\
	babička na mě šetří v Panamě\\
	a já si tu žiju náramě!\\
\end{verse}



\end{song}

\end{minipage}
\end{document}
